%% 由 build/emit_tex.py 生成（生成物，勿手改）；定制请放 user_style.tex / user_meta.yaml
\chapter{习题2 参考答案}\label{ux4e60ux98982-ux53c2ux8003ux7b54ux6848}

\begin{quote}
说明：第2章（解析函数）主要习题与参考答案（经典教材整理）。公式请用
\(...\) / \[...\] 阅读。
\end{quote}

\section{习题}\label{ux4e60ux9898-1}

\begin{enumerate}
\def\labelenumi{\arabic{enumi}.}
\item
  利用导数定义推出：

  \begin{enumerate}
  \def\labelenumii{(\arabic{enumii})}
  \tightlist
  \item
    \((z^n)'=nz^{n-1}\)（\(n\) 为正整数）；
  \item
    \(\left(\dfrac{1}{z}\right)'=-\dfrac{1}{z^2}\)。
  \end{enumerate}
\item
  下列函数何处可导？何处解析？

  \begin{enumerate}
  \def\labelenumii{(\arabic{enumii})}
  \tightlist
  \item
    \(f(z)=x-iy^2\)；
  \item
    \(f(z)=2x^3+3iy^3\)；
  \item
    \(f(z)=xy^2+ix^2y\)；
  \item
    \(f(z)=\sin x\cosh y+i\cos x\sinh y\)。
  \end{enumerate}
\item
  指出下列函数的解析性区域，并求其导数：

  \begin{enumerate}
  \def\labelenumii{(\arabic{enumii})}
  \tightlist
  \item
    \((z-1)^5\)； (2) \(z^3+2iz\)； (3) \(\dfrac{1}{z^2-1}\)； (4)
    \(\dfrac{az+b}{cz+d}\)（\(ad-bc\neq0\)）。
  \end{enumerate}
\item
  复变函数的可导性与解析性有什么区别？判断函数解析性有哪些方法？
\item
  判断下述命题的真假，并举例说明：

  \begin{enumerate}
  \def\labelenumii{(\arabic{enumii})}
  \tightlist
  \item
    若 \(f(z)\) 在 \(z_0\) 连续，则 \(f'(z_0)\) 存在；
  \item
    若 \(f'(z_0)\) 存在，则 \(f(z)\) 在 \(z_0\) 解析；
  \item
    若 \(z_0\) 是 \(f(z)\) 的奇点，则 \(f(z)\) 在 \(z_0\) 不可导；
  \item
    若 \(z_0\) 是 \(f(z)\) 和 \(g(z)\) 的奇点，则 \(z_0\) 也是
    \(f(z)+g(z)\) 的奇点；
  \item
    若 \(u(x,y),v(x,y)\) 的偏导数存在，则 \(f(z)=u+iv\) 可导；
  \item
    若 \(f(z)=u+iv\) 在区域 \(D\) 内解析，且 \(u\) 为实常数，则 \(f(z)\)
    在 \(D\) 内为常数（\(v\) 为实常数时同理）。
  \end{enumerate}
\item
  求下列各值：

  \begin{enumerate}
  \def\labelenumii{(\arabic{enumii})}
  \tightlist
  \item
    \(e^{2+i}\)； (2) \(\operatorname{Ln}(3+4i)\)； (3) \(\ln 1\)； (4)
    \((1+i)^i\)； (5) \(\sin i\)。
  \end{enumerate}
\item
  若 \(z=x+iy\)，证明：

  \begin{enumerate}
  \def\labelenumii{(\arabic{enumii})}
  \tightlist
  \item
    \(\sin z=\sin x\cosh y+i\cos x\sinh y\)；
  \item
    \(|\sin z|^2=\sin^2 x+\sinh^2 y\)；
  \item
    \(|\cos z|^2=\cos^2 x+\sinh^2 y\)。
  \end{enumerate}
\item
  证明：当 \(y\to\infty\) 时，\(|\sin(x+iy)|\) 与 \(|\cos(x+iy)|\)
  都趋于无穷大。
\end{enumerate}

\section{参考答案}\label{ux53c2ux8003ux7b54ux6848-1}

\begin{enumerate}
\def\labelenumi{\arabic{enumi}.}
\item
  由定义
  \(\lim_{\Delta z\to0}\dfrac{(z+\Delta z)^n-z^n}{\Delta z}=nz^{n-1}\)；
  对
  \(1/z\)：\(\lim_{\Delta z\to0}\dfrac{1/(z+\Delta z)-1/z}{\Delta z}=\lim_{\Delta z\to0}\dfrac{-1}{z(z+\Delta z)}=-\dfrac1{z^2}\)。
\item
  用 C-R 方程判断：

  \begin{enumerate}
  \def\labelenumii{(\arabic{enumii})}
  \tightlist
  \item
    \(u=x,\ v=-y^2\)，\(u_x=1,v_y=-2y\)，\(u_y=0,v_x=0\)，仅当
    \(1=-2y\)，即 \(y=-1/2\) 时可导，处处不解析；
  \item
    \(u=2x^3,\ v=3y^3\)，满足 C-R 仅当 \(6x^2=9y^2\)，即在直线
    \(2x^2=3y^2\) 上可导，处处不解析；
  \item
    \(u=xy^2,\ v=x^2y\)，满足 C-R 仅当 \(y^2=x^2\) 且 \(2xy=-2xy\)，即
    \(z=0\) 可导，处处不解析；
  \item
    \(u=\sin x\cosh y,\ v=\cos x\sinh y\)，全平面满足 C-R，故
    \(f(z)=\sin z\) 在全平面解析。
  \end{enumerate}
\item
  \begin{enumerate}
  \def\labelenumii{(\arabic{enumii})}
  \tightlist
  \item
    \(f'(z)=5(z-1)^4\)，全平面解析；
  \item
    \(f'(z)=3z^2+2i\)，全平面解析；
  \item
    除 \(z=\pm1\) 外解析，\(f'(z)=\dfrac{-2z}{(z^2-1)^2}\)，\(z=\pm1\)
    为奇点；
  \item
    除 \(z=-d/c\) 外解析，\(f'(z)=\dfrac{ad-bc}{(cz+d)^2}\)。
  \end{enumerate}
\item
  可导是局部性质（一点处），解析是``邻域内处处可导''（区域性质）：一点可导不一定解析。判断方法：①
  定义；② C-R 方程 + 实部虚部可微（充要条件）。
\item
  \begin{enumerate}
  \def\labelenumii{(\arabic{enumii})}
  \tightlist
  \item
    假：\(f(z)=|z|^2\) 全平面连续，但只在 \(z=0\) 可导；
  \item
    假：\(f(z)=|z|^2\) 在 \(z=0\) 可导但该点不解析；
  \item
    假：如 \(f(z)=\bar z\)，处处连续且处处不可导，但无``孤立奇点''；
  \item
    假：如 \(f(z)=g(z)=\operatorname{Re}z\) 有公共``不可导点''，但
    \(f+g\)
    仍处处不可导？更恰当举例：\(f(z)=\dfrac1z\)，\(g(z)=-\dfrac1z+\dfrac1{z-1}\)
    等；结论是命题不真；
  \item
    假：偏导数存在不足以保证 C-R 满足；
  \item
    真：\(u\) 为常数时，由 C-R 方程 \(v_x=-u_y=0,\ v_y=u_x=0\)，所以
    \(v\) 也为常数，\(f\) 为常数。
  \end{enumerate}
\item
  \begin{enumerate}
  \def\labelenumii{(\arabic{enumii})}
  \tightlist
  \item
    \(e^{2+i}=e^2(\cos1+i\sin1)\)；
  \item
    \(\operatorname{Ln}(3+4i)=\ln5+i(\arctan\frac43+2k\pi)\)；
  \item
    \(\ln1=\ln|1|+i\arg1=0\)（主值）；
  \item
    \((1+i)^i=e^{i\operatorname{Ln}(1+i)}=e^{i(\ln\sqrt2+i\pi/4)}=e^{-\pi/4}(\cos\ln\sqrt2+i\sin\ln\sqrt2)\)；
  \item
    \(\sin i=i\sinh1\)。
  \end{enumerate}
\item
  \begin{enumerate}
  \def\labelenumii{(\arabic{enumii})}
  \tightlist
  \item
    由 \(e^{i(x+iy)}\) 与 \(e^{-i(x+iy)}\) 定义直接展开； (2)(3) 用 (1)
    及 \(\cosh^2y-\sinh^2y=1\) 化简。
  \end{enumerate}
\item
  \(|\sin(x+iy)|^2=\sin^2x+\sinh^2y\)，当 \(|y|\to\infty\) 时
  \(\sinh^2y\to\infty\)；\(\cos\) 同理。
\end{enumerate}
